\subsection{Setting and quantities}
A \emph{system} has a finite set of candidate configurations $\mathcal C=\{1,\dots,K\}$ (prompt wordings, demonstration sets, option layouts). A configuration $c$ has an unknown population accuracy $\mu_c$ on the benchmark's item distribution. A \emph{tuning set} $T$ of $n$ items is drawn, a \emph{search procedure} $\mathcal S$ sees only the correctness of the candidates on $T$ and returns a configuration $\hat c$, and the \emph{naive report} is the accuracy of $\hat c$ on the same items, $R_{\text{naive}}=\mathrm{acc}_T(\hat c)$. The target of any honest report is $\mu_{\hat c}$, the accuracy of the \emph{selected} configuration on fresh items. We call
$\Delta = R_{\text{naive}}-\mu_{\hat c}$ the \emph{optimism gap}. It has three sources: selection bias (the maximum of $K$ noisy scores), extra adaptivity of the search (rounds, budget allocation), and leakage when demonstrations are taken from $T$. We study the first two and, in the synthetic setting, leakage.

The unit of repetition throughout is the \emph{whole pipeline}: draw $T$, search, select, report, compare with the truth. For a reporting method $M$ with point report $\hat\theta_M$ and optional 95\% interval, over $R$ repetitions we measure bias $\mathbb E[\hat\theta_M-\mu_{\hat c}]$, RMSE, interval coverage of $\mu_{\hat c}$, interval width, and the number of items consumed. Every metric carries a Monte-Carlo standard error, and coverage is compared with the binomial Monte-Carlo band around $0.95$, not with $0.95$ itself.

\subsection{Search procedures}
All procedures see an array of $K\times n$ correctness indicators. \emph{Best-of-$K$} (full evaluation of every candidate on $T$, argmax with random tie-breaking) models grid or random search. \emph{Successive halving} splits $T$ into $\lceil\log_2K\rceil$ disjoint blocks and keeps the better half of the survivors each round; \emph{successive rejects} is a fixed-budget best-arm identification rule (the framing of prompt selection used by \citet{shi2024triple}), run here with a budget of $nK/4$ evaluations; we implement it ourselves and do not run the algorithm of that paper. \emph{Greedy coordinate ascent} over a factorial grid imitates multi-round instruction/demonstration search. In every case the naive report is the winner's score on the evaluations the procedure collected for it.

\subsection{Reporting methods}
\textbf{Naive} (with Wilson interval). \textbf{Held-out split}: select on half of $T$, report on the other half; unbiased for $\mu_{\hat c}$ but it spends half the items and ships a configuration chosen from $n/2$ items. \textbf{Fresh set}: select on $n_{\text{tune}}$ items and report the winner on $m$ new items. \textbf{$k$-fold cross-fitting} ($k=5$): run the whole procedure on $k-1$ folds and score its pick on the held-out fold; its estimand is the performance of the \emph{procedure} on $(1-1/k)n$ items. \textbf{Bootstrap optimism correction}: resample $T$, rerun the procedure, and subtract the mean difference between the score on the resample and on $T$; we also use the out-of-bag and the .632 variants. \textbf{Simultaneous bound}: an interval whose critical value is the $(1-\alpha)$ quantile of the bootstrap maximum of $|\hat p^*_k-\hat p_k|/\mathrm{se}_k$ over all $K$ candidates, valid for any data-dependent selection \citep{chernozhukov2013gaussian,berk2013valid}; its point report remains the naive score. \textbf{Conditional (winner-aware) inference} for the full-evaluation argmax: under a Gaussian approximation of the $K$ scores the selection event is a polyhedron, and inverting the truncated-normal law of the winner's score gives a median-unbiased estimate and an equal-tailed interval, and a hybrid with a projection interval; this follows \citet{andrews2024winners} but is our own implementation. \textbf{Ladder-type thresholded report} \citep{blum2015ladder}: for sequential submission only, the reported best-so-far score moves only when a candidate beats it by more than $1/\sqrt n$.

\subsection{Gap prediction}
For best-of-$K$ with Gaussian tuning scores of covariance $\Sigma/n$ (the question-level covariance of the candidates' correctness), the expected gap is $\mathbb E[\epsilon_{\hat c}]$, with $\epsilon\sim N(0,\Sigma/n)$ and $\hat c=\arg\max_c(\mu_c+\epsilon_c)$. For independent noise with standard deviation $s=\sqrt{p(1-p)/n}$ and true-accuracy spread $\tau$ this is approximately $\frac{s^2}{\sqrt{s^2+\tau^2}}\,\mathbb E[\max_{k\le K}Z_k]$ with $Z_k$ standard normal. We use these as empirical approximations to compare with simulation, not as theorems.

\subsection{Synthetic and real settings}
\paragraph{Synthetic (exact truth).} A Gaussian-copula item pool: configuration accuracies $p_c=m+\tau z_c$ with $z_c$ standard normal, an item latent shared across configurations with correlation $\rho$, and correctness when the latent falls below $\Phi^{-1}(p_c)$. The truth of the selected configuration is its mean over the pool ($N=20{,}000$ items, $128$ configurations; Appendix~\ref{app:defs}). Optional arms add generation noise (each evaluation flipped with probability $\varepsilon$), leakage (a number of tuning items shown with their answers in the prompt) and one outlier configuration.
\paragraph{Real (small LLM, MMLU-Pro).} The LLM is the system under tuning: Qwen2.5-0.5B-Instruct (fp32, CPU, no sampling) scores a fixed pool of $2{,}000$ MMLU-Pro test items (revision \texttt{b189ec7}, full hash in the reproduction notes) under a pool of prompt configurations that varies the instruction wording, the demonstration set and the option layout ($60$ configurations), by reading the log-probabilities of the ten option letters at the answer position. Scoring is a single deterministic forward pass, so the correctness matrix is a fixed array and any search can be replayed on it. Tuning sets are drawn from the pool, and the truth is the accuracy of the selected configuration on the pool items not in $T$ (\emph{disjoint} arm; an i.i.d. arm against the pool mean is a check). All real statements are conditional on this pool, model and scoring protocol.
